\section{Background and Motivation}
\subsection{Four events that need not coincide}
For a producer $u$ and consumer $v$, distinguish producer success $S_u$,
consumer dispatch $D_v$, consumer input readiness $R_v$, and consumer execution
$E_v$. A valid execution requires the relevant inputs at $E_v$; it does not
inherently require those inputs to be present on the selected worker at $D_v$.
For an intermediate with asynchronous backup, denote backup admission by $B_u$.
A policy may permit $S_u \prec D_v$ and $R_v \prec E_v$ without imposing
$B_u \prec D_v$. This is an ordering permission, not a guarantee that a
particular transfer or backup finishes before another event.

Requested outputs introduce a different boundary. A user observing a final
result needs the service's promised durability, so physical task completion
alone cannot establish workflow-result availability. Separating these events
allows the runtime to move successors forward while retaining a stronger
condition at the user-visible result interface.

The bottleneck is observable before introducing a new scheduling policy.
Figure~\ref{fig:bottleneck-motivation}(a) increases one-MiB durable objects
from 64 to 1,024 at a fixed four-worker topology: service time grows from
0.87 to 13.64 seconds while manager owner work stays below 28 ms after
isolation. In (b), 4 GiB remote fan-in throughput falls from 542.8 to
442.9 MiB/s as endpoints grow from 16 to 64, with Controller CPU near
36--37 seconds. The matched 20,000-task comparison (c)--(d) sends 64.1 MiB
through the coupled manager and completes in 160.4 seconds, versus 0.04 MiB
and 10.0 seconds for the separated path (16.1$\times$ median). These results
motivate separating logical scheduling from data-lifecycle service.

\begin{figure*}[t]
\centering
\includegraphics[width=.98\textwidth]{bottleneck-motivation.pdf}
\caption{Motivation evidence from real runs. (a) Fixed four-worker by two-core
scale-up with one-MiB durable outputs; the data-plane service interval grows
with objects while manager owner work stays small after separation. (b) Remote
Controller fan-in with 4 GiB total payload; adding endpoints exposes a shared
service boundary. (c)--(d) Matched 20,000-task data-intensive comparison: the
coupled path carries manager data-plane traffic and is slower, while the
separated path keeps payload service outside the scheduler. Bars use retained
medians; source hashes and raw runs are in the artifact.}
\label{fig:bottleneck-motivation}
\end{figure*}

\subsection{Queue depth is not executable capacity}
Let $Q$ be the maximum number of assigned function calls on a worker, and $W$
its current execution window. For $C$ allocated cores, a queue with $Q>C$
can buffer descriptors without requiring $Q$ executing processes or $Q$ fully
prepared input sandboxes. A conventional single limit obscures which form of
concurrency is increasing. DataVine makes $Q$ and $W$ distinct and uses the
availability of an execution opportunity to gate preparation.

A simple operational decomposition helps interpret measurements. A task's
observed worker residence includes waiting before preparation, preparation
itself, waiting after preparation, and execution. Increasing $Q$ primarily
changes what work is available locally; increasing $W$ changes competition
among executing or input-waiting calls. Neither necessarily reduces the
critical path. On a computation-bound phase, more processes may only increase
contention. On a data-waiting phase, a larger window may expose otherwise idle
CPU capacity. Measurements must therefore cover both regimes and transitions.

\subsection{What the motivating comparison must preserve}
A comparison against a baseline that sends every intermediate through a central
client would not isolate DataVine's control boundary. Our TaskVine baseline uses
its futures interface, temporary intermediate outputs, fork-based function
libraries, and the same kernel code. It benefits from the existing runtime's
peer-transfer and caching machinery. Every graph retains individual tasks;
there is no task fusion or replacing a chain with a single process.

Task grouping also advances some dependency knowledge toward workers~\cite{grouping}.
It must not be dismissed as fusion. Similarly, workflow-aware prefetch can
overlap movement and computation even if assignment requires prepared inputs.
The relevant distinction is the event at which assignment becomes legal, not
whether a system has asynchronous operations. Our eager/deferred ablation
isolates preparation timing within DataVine; it is not presented as a complete
reimplementation of either grouping or a prepared-node scheduler.
